\documentclass[11pt,a4paper]{article}

\usepackage[T1]{fontenc}
\usepackage[utf8]{inputenc}
\usepackage{lmodern}
\usepackage[margin=1in]{geometry}
\usepackage{hyperref}
\usepackage{booktabs}
\usepackage{graphicx}
\usepackage{xcolor}
\usepackage{url}
\usepackage{microtype}
\usepackage{authblk}

\hypersetup{
  colorlinks=true,
  linkcolor=black,
  urlcolor=blue!60!black,
  citecolor=blue!60!black,
}

\title{\textbf{Japan Travel MCP: A 17-Language Open Dataset and Model Context
Protocol Server for Japanese Tourism Information}}

\author[1]{KJ Sunada\thanks{\texttt{kjsunada@huggingface.co}}}
\affil[1]{KabuK Style Inc., Tokyo, Japan}

\date{April 2026}

\begin{document}
\maketitle

\begin{abstract}
We introduce \textbf{Japan Travel MCP}, an open dataset and reference Model
Context Protocol (MCP) server that exposes structured, multilingual Japanese
tourism information for AI agents. The dataset covers all 47 prefectures and
1{,}938 local-government entities of Japan, providing 13{,}394 tourist-spot
descriptions in 17 languages (227{,}698 description cells), 13{,}961 canonical
multilingual names (237{,}337 name cells), 40{,}128 unified accommodation
records, and 690 officially-designated cultural records spanning the
Ministry of Agriculture, Forestry and Fisheries' Geographical Indications,
the Ministry of Economy, Trade and Industry's Traditional Crafts, the Agency
for Cultural Affairs' Japan Heritage stories and Important Intangible
Cultural Properties, and Japan's UNESCO Intangible Cultural Heritage
inscriptions. All data is built exclusively from authoritative public
sources (municipal websites, Wikidata, OpenStreetMap, official designation
registries) under a CC BY 4.0 license, and is published as a
ready-to-consume Hugging Face dataset accompanied by a TypeScript MCP
server published to npm. We describe the data pipeline, the
glossary-grounded multilingual translation methodology using Anthropic's
Claude Sonnet 4.6 batch API, the entity-matching algorithm that unifies
hotel records across Wikidata and OpenStreetMap, and the design choices
behind the MCP tool surface. To our knowledge, this is the first
publicly-available, multilingual, license-clear, AI-agent-ready resource for
Japanese tourism that prioritises the long tail of municipalities outside
Tokyo and Kyoto.

\medskip\noindent\textbf{Keywords:} multilingual datasets, tourism information,
model context protocol, low-resource languages, Japanese NLP, open data,
entity matching.
\end{abstract}

\section{Introduction}
\label{sec:intro}

Japan's tourism information is published across thousands of municipal,
prefectural, and ministerial websites, primarily in Japanese. While Tokyo
and Kyoto enjoy comprehensive coverage in commercial tourism platforms,
the long tail of Japan's 1{,}938 local-government entities is poorly
represented in the AI-accessible knowledge graph. Large language models
asked about tourism in, say, Tottori, Kochi, or Saga frequently default to
generic, stale, or non-existent information. This is not a model failure
\textit{per se}: it is a data-availability failure. The information exists,
but in forms that are neither machine-readable, nor multilingual, nor
under licenses that permit redistribution.

We introduce \textbf{Japan Travel MCP}, an open project that closes this
gap on three fronts simultaneously:

\begin{enumerate}
  \item A \textbf{multilingual dataset} of Japanese tourism information,
  covering all 47 prefectures and 1{,}938 local-government entities, in
  17 languages, distributed under CC BY 4.0 with full source attribution.
  \item A \textbf{reference MCP server} that exposes the dataset to AI
  agents through ten tools, distributable via \texttt{npm} for stdio MCP
  clients (Claude Desktop, Cursor, Windsurf) and via Docker for HTTP-MCP
  hosts (Hugging Face Spaces).
  \item A \textbf{transparent build pipeline}: every record carries the
  authority that designated it, the URL it was scraped from, and the date
  of last refresh. The pipeline runs daily on a public GitHub Actions
  schedule.
\end{enumerate}

This paper describes the dataset, the pipeline, and the reference
application. Section~\ref{sec:data} details the dataset and source
hierarchy; Section~\ref{sec:method} describes the data pipeline including
the multilingual translation methodology; Section~\ref{sec:coverage}
reports coverage statistics; Section~\ref{sec:app} describes the MCP
reference application; Section~\ref{sec:limitations} discusses limitations
and open problems; Section~\ref{sec:related} surveys related work.

\section{Dataset}
\label{sec:data}

\subsection{Sources and licensing}

The dataset is composed of seven layers, each derived from publicly-available
sources under licenses that permit redistribution with attribution:

\begin{itemize}
  \item \textbf{Layer 1 — Municipal tourism pages.} Tourism content scraped
  from the official websites of all 1{,}741 municipalities and 197
  designated-city wards (1{,}938 entities total). Robots.txt is respected
  on every domain. The scraper enforces a minimum 5-second per-domain
  interval and a 30-day rolling refresh; this is slower than Googlebot,
  by design, to avoid imposing load on small municipal servers.

  \item \textbf{Layer 2 — Wikidata attractions.} 41{,}404 ja-anchored
  tourist-attraction entities (\texttt{Q570116} and its subclasses),
  obtained via SPARQL queries to query.wikidata.org. Wikidata content is
  CC0; we record the QID and Wikidata URL on every record.

  \item \textbf{Layer 3 — Hotel master.} 40{,}128 accommodation records
  unified from Wikidata (CC0) and OpenStreetMap (ODbL). Section~\ref{sec:matching}
  describes the matching algorithm.

  \item \textbf{Layer 4 — Wikipedia sitelinks.} For every Wikidata entity
  with a Japanese-Wikipedia anchor, we collect parallel sitelinks across
  the 17 target languages, forming a high-quality canonical-name backbone
  for the multilingual layer.

  \item \textbf{Layer 5 — Multilingual translation layer.} 17-language
  names (237{,}337 cells) and 200--300-character tourism descriptions
  (227{,}698 cells), produced by a Claude Sonnet 4.6 batch pipeline
  grounded in a project-wide canonical glossary
  (Section~\ref{sec:translation}).

  \item \textbf{Layer 6 — Official designation registries.} 690 records
  drawn directly from five Japanese government designation systems:

  \begin{itemize}
    \item Ministry of Agriculture, Forestry and Fisheries (MAFF)
    Geographical Indications (172 records);
    \item Ministry of Economy, Trade and Industry (METI) Traditional
    Crafts (Dent\=o K\=ogeihin, 231 records);
    \item Agency for Cultural Affairs (Bunka-ch\=o) Japan Heritage
    (Nihon Isan, 104 stories);
    \item Bunka-ch\=o Important Intangible Cultural Properties and
    Important Intangible Folk Cultural Properties (125 records, mirrored
    via Wikidata);
    \item UNESCO Intangible Cultural Heritage inscriptions for Japan
    (58 records, mirrored via Wikidata).
  \end{itemize}

  All Layer 6 records are translated into the same 17 languages
  (11{,}730 cells).

  \item \textbf{Layer 7 — Operational metadata.} Per-municipality
  scrape-state, refresh timestamps, and resolved official-URL mappings,
  which support reproducibility and incremental maintenance.
\end{itemize}

\subsection{Source-selection principle}

We adopt a strict ``official build-up'' policy: only records that an
authoritative public body has designated, and that we can scrape directly
from that body's own publication, are admitted. We deliberately exclude
editorial picks, AI-curated lists, user-generated content, OTA scraping
(which would violate ToS), and any source requiring authentication.
This policy sacrifices breadth for verifiability: a tourism record in
Japan Travel MCP can always be traced back to a public Japanese
government, prefectural, or municipal website.

\section{Methodology}
\label{sec:method}

\subsection{Municipal scraping pipeline}

For each of the 1{,}938 local-government entities, the pipeline:
\begin{enumerate}
  \item resolves the official tourism URL, with manual overrides for
  edge cases (defunct subdomains, language-toggle redirects);
  \item validates robots.txt, and aborts the domain if the
  tourism path is disallowed;
  \item performs a tourism-keyword breadth-first crawl, extracting
  pages whose title or anchor text matches a curated set of Japanese
  and English tourism keywords (e.g.\ \emph{kanko}, \emph{sightseeing},
  \emph{festival}, \emph{\d{}观光});
  \item parses each page through a heuristic extractor producing
  \texttt{(name, description, address, coordinates, language)} tuples;
  \item canonicalises spot identifiers via SHA-256 of the canonical URL,
  preserving stability across re-scrapes.
\end{enumerate}

The pipeline runs on a daily GitHub Actions cron (03:00 JST) processing
70 entities per day, ensuring every entity is refreshed within 28 days.

\subsection{Hotel entity matching}
\label{sec:matching}

The hotel master is built by merging Wikidata and OpenStreetMap. Two
records are considered the same property if they fall within 100 meters
of each other and share a sufficiently similar name (a hybrid metric
over kanji, kana, and romaji forms). Confirmed clusters receive
\texttt{confidence: confirmed}; singleton records are retained as
\texttt{confidence: singleton}; ambiguous matches (similar names with
borderline geographic overlap) are written to a public \texttt{review/}
queue for community resolution via pull request.

\subsection{Multilingual translation}
\label{sec:translation}

The translation layer employs a three-tier source hierarchy:

\begin{enumerate}
  \item \textbf{Wikipedia article titles} (via Wikidata sitelinks):
  human-curated, peer-reviewed translations are used directly when
  available.
  \item \textbf{Project canonical glossary}: a hand-maintained term
  list for tourism-specific suffix and proper-noun rules
  (\emph{jinja} → Shrine, \emph{tera} → Temple, \emph{j\=o} → Castle;
  modified Hepburn romanization), loaded as a cached system prompt for
  AI-translation consistency.
  \item \textbf{Claude Sonnet 4.6 batch translation}: fills gaps that
  Wikipedia does not cover. We use the Anthropic Batch API for the
  50\% cost discount; 13{,}961 names + 13{,}394 descriptions + 690
  designation records were generated for approximately USD~111 total.
\end{enumerate}

The translation prompt is constrained to enforce JSON-only output with
language-natural quotation marks (\texttt{「」} for ja/zh, \texttt{« »}
for fr/ru/it/pt/es, \texttt{„''} for de) to prevent JSON-string corruption
from improperly-escaped quotes — a failure mode that affected approximately
33\% of records in our initial run and required a key-vs-value-aware
JSON repair function in post-processing.

\section{Coverage}
\label{sec:coverage}

\begin{table}[h]
\centering
\begin{tabular}{lrr}
\toprule
\textbf{Layer} & \textbf{Records} & \textbf{Cells} \\
\midrule
Multilingual canonical names                   & 13{,}961 & 237{,}337 \\
Multilingual tourism descriptions (200--300 ch)& 13{,}394 & 227{,}698 \\
Official-designation translations              &     690  &  11{,}730 \\
Wikidata attraction backbone (raw)             & 41{,}404 &        — \\
Hotel master (Wikidata + OSM merged)           & 40{,}128 &        — \\
Municipalities + designated-city wards covered &  1{,}938 &        — \\
Prefectures covered                            &     47   &        — \\
\bottomrule
\end{tabular}
\caption{Dataset coverage as of 2026-04-30.}
\label{tab:coverage}
\end{table}

Coverage is uniform across the 17 target languages: every entity in the
description-translation set has descriptions in all 17 languages by
construction. Per-prefecture entity counts vary substantially, with
Hokkaido, Kyoto, and Tokyo carrying the largest absolute counts and
several rural prefectures (Yamanashi, Tokushima, Saga) holding fewer
than 200 entities. The point of the dataset, however, is that
\textit{all} 47 prefectures are represented at non-trivial depth — there
is no per-prefecture dropout below 100 entities.

\section{MCP Reference Application}
\label{sec:app}

The reference MCP server exposes ten tools through both stdio and
Streamable HTTP transports (the latter targeting hosted deployments
such as Hugging Face Spaces). The tool surface is summarised in
Table~\ref{tab:tools}.

\begin{table}[h]
\centering
\small
\begin{tabular}{ll}
\toprule
\textbf{Tool} & \textbf{Purpose} \\
\midrule
\texttt{search\_area}        & Cross-language name/keyword search \\
\texttt{get\_spots}          & Tourist spots by prefecture / municipality \\
\texttt{get\_hotels}         & Accommodations by area or geographic radius \\
\texttt{get\_transport}      & Spot location + official access URL \\
\texttt{get\_events}         & Festivals (Wikidata SPARQL, in-memory cache) \\
\texttt{get\_multilingual}   & Lightweight multilingual name lookup \\
\texttt{get\_description}    & 17-language tourism description (signature) \\
\texttt{get\_local\_specialty}    & MAFF GI + METI Traditional Crafts \\
\texttt{get\_traditional\_arts}   & Bunka-ch\=o Intangible CP + UNESCO ICH \\
\texttt{get\_japan\_heritage}     & Bunka-ch\=o Japan Heritage stories \\
\bottomrule
\end{tabular}
\caption{Tool surface of the reference MCP server.}
\label{tab:tools}
\end{table}

The npm package (\texttt{japan-travel-mcp}) is approximately 386~kB
compressed and downloads the runtime data (approximately 270~MB) from
the Hugging Face dataset on first run, caching to
\texttt{\textasciitilde{}/.japan-travel-mcp/data/}.

\section{Limitations and Open Problems}
\label{sec:limitations}

\paragraph{Entity matching.} The hotel matcher uses a 100-meter geographic
threshold combined with a name-similarity metric. Properties with similar
names in close proximity (e.g.\ a hotel and an attached restaurant) can
produce false-positive merges; the manual \texttt{review/} queue is the
current mitigation but does not scale to higher data volumes.

\paragraph{Low-resource languages.} For Hindi, Tagalog, and Thai we have
no Wikipedia sitelinks for many entities, and translation quality therefore
relies entirely on Claude Sonnet 4.6 with glossary grounding. We report
self-rated confidence (\texttt{high}/\texttt{medium}/\texttt{low}) per
record but lack a held-out evaluation against native speakers; this is
priority future work.

\paragraph{Freshness model.} The 30-day rolling refresh is appropriate for
near-static municipal tourism content but is too coarse for time-sensitive
information such as festival dates and seasonal closures. The
\texttt{get\_events} tool currently performs a live SPARQL query at
request time as a partial mitigation.

\paragraph{Coverage of designation registries.} The 125 records under the
Bunka-ch\=o Important Intangible Cultural Property layer are sourced via
Wikidata, which mirrors only a subset of the full bunka.go.jp registry.
Direct ingestion from the official database is hindered by the absence of
a public CSV / API export; this is an open community task.

\section{Related Work}
\label{sec:related}

Existing tourism datasets either (a) cover Japan in a single language
(JNTO open data; mostly English), (b) lack the entity-level granularity
required for AI agents (Wikipedia / Wikidata alone, which are sparse
outside major cities), or (c) carry restrictive licenses that prevent
redistribution (commercial APIs such as Google Places, Foursquare).

The Model Context Protocol \cite{mcp-spec}, introduced by Anthropic in
late 2024, has rapidly become the de-facto interface between AI agents
and external tools. To our knowledge, Japan Travel MCP is the first
country-scale, multilingual, license-clear MCP server for tourism
information.

Multilingual canonical-glossary-grounded translation has antecedents in
machine-translation research \cite{lample2018phrase}, but its application
to large-scale tourism content with explicit suffix-rule consistency is,
to our knowledge, novel.

\section{Availability}
\label{sec:availability}

\begin{itemize}
  \item \textbf{Code:} \url{https://github.com/ookami0210/japan-travel-mcp} (MIT license)
  \item \textbf{Dataset:} \url{https://huggingface.co/datasets/kjsunada/japan-travel-mcp-data} (CC BY 4.0)
  \item \textbf{npm package:} \url{https://www.npmjs.com/package/japan-travel-mcp}
\end{itemize}

The full data pipeline, including scrapers, translators, and the MCP
server itself, is published under permissive open licenses to enable
reproducibility, extension, and adaptation to other countries with
similar information-fragmentation patterns.

\section*{Acknowledgements}

We thank the maintainers of Wikidata, OpenStreetMap, and the Japanese
municipal tourism websites whose work this dataset compiles into a
unified, multilingual form. We thank Anthropic for the Model Context
Protocol specification and the Claude Sonnet 4.6 Batch API. The first
external code review and test infrastructure was contributed by
Russlan Iablonovskii.

\begin{thebibliography}{9}

\bibitem{mcp-spec}
Anthropic.
\textit{Model Context Protocol Specification}.
2024. \url{https://modelcontextprotocol.io/}

\bibitem{lample2018phrase}
G.~Lample, M.~Ott, A.~Conneau, L.~Denoyer, M.A.~Ranzato.
Phrase-based \& neural unsupervised machine translation.
In \textit{Proceedings of EMNLP}, 2018.

\end{thebibliography}

\end{document}
